%======================================================================
% 附录J：经典教材逐章对照与延伸阅读
%======================================================================
% 章标题 \appchapter{附录J：...} 由主文件发出，本文件只提供正文。

附录 E 回答“该买哪本书”，本附录回答一个更具体的问题。\textbf{本书某一章卡住了，去那本书的哪一部分看？}

下面把本书每一章对应到 Chen 与 Fitzpatrick 的章，以及更专门的专著里由谁讲得最清楚。

\begin{warning}{关于章号的使用方式}{}
不同版次的章号会变。本附录的章号取自本节末列出的具体版次（Chen 第 3 版；Fitzpatrick 第 2 版），并已在核对日期当日逐条比对目录。若你手上是别的版次，请按\textbf{主题名}而不是章号查找；本附录每一条都同时给了主题名。
\end{warning}


\section{与附录 E 的分工}

\begin{itemize}
\item \textbf{附录 E}：选书。按水平分级，每条给出“解决本书哪个疑问”。
\item \textbf{本附录}：定位。按下表逐章对照，给出“这一章在哪本书里讲得更透、以及它们讲法上的差别”。
\item \textbf{附录 K}：练手。给出八道跨教材经典例题，并注明出处。
\end{itemize}


\section{七本书的“性格”速写}

\begin{enumerate}
\item \textbf{Chen，\textit{Introduction to Plasma Physics and Controlled Fusion}, 3rd ed.}\\
语言最通俗，物理图像优先，每章有大量数值估计题。适合“先建立直觉”。它的第 9、10 章（Special Plasmas、Plasma Applications）是第 3 版新增内容。与本书第 10--12 章互补。

\item \textbf{Fitzpatrick，\textit{Plasma Physics: An Introduction}, 2nd ed.}\\
推导最完整，且\textbf{全书 HTML 免费在线}。它的第 4 章把流体闭合（Chapman--Enskog、Braginskii）讲得比本书更形式化。第 9 章专门讲磁重联，包括撕裂模的渐近匹配与 Sweet--Parker 模型；本书受篇幅限制只给了量级估计。适合“要把推导补齐”。

\item \textbf{Goldston \& Rutherford，\textit{Introduction to Plasma Physics} (IOP, 1995)。}\\
PPPL 的研究生教材，数值例子贴近聚变实验装置。定位按\textbf{主题}给出（单粒子运动、等离子体流体、碰撞过程、波、动理学），不依赖具体章号。

\item \textbf{Wesson，\textit{Tokamaks} (Oxford University Press)。}\\
托卡马克物理的百科全书式专著。当你想知道“书上这个不稳定性在真实装置上的增长率与阈值是多少”，查它。

\item \textbf{Krall \& Trivelpiece，\textit{Principles of Plasma Physics} (McGraw-Hill, 1973)。}\\
经典“大书”，动理学与波的数学处理最完整，适合当作长期案头参考。

\item \textbf{Hazeltine \& Meiss，\textit{Plasma Confinement} (Addison-Wesley, 1992)。}\\
从约束与输运的角度组织内容，对回旋动理学与输运理论的直觉解释出色。

\item \textbf{Freidberg，\textit{Ideal MHD} (Cambridge University Press, 2014)。}\\
理想 MHD 平衡与稳定性的现代标准参考，能量原理（$\delta W$）讲得最系统；配合本书第 5、9 章使用。
\end{enumerate}


\section{逐章对照表}

\small
\begin{longtable}{@{}p{2.7cm}p{2.5cm}p{2.5cm}p{3.2cm}p{3.0cm}@{}}
\caption{本书各章与经典教材的对照（Fitz 指 Fitzpatrick 第 2 版，Chen 指 Chen 第 3 版）}\label{tab:app-j-map}\\
\toprule
本书 & 主题 & Chen & Fitz & 其他专著 \\
\midrule
\endfirsthead
\toprule
本书 & 主题 & Chen & Fitz & 其他专著 \\
\midrule
\endhead
\midrule
\multicolumn{5}{r}{\small（续下页）}\\
\endfoot
\bottomrule
\endlastfoot
第 1 章 & 数学与统计力学预备知识 & 第 1 章、附录 A（单位、矢量关系） & 无专章，工具散见各章 & Krall \& Trivelpiece 的数学附录 \\
第 2 章 & 什么是等离子体 & 第 1 章 & 第 1 章（判据与参数） & --- \\
第 3 章 & 单粒子运动 & 第 2 章 & 第 2 章（平均化、绝热不变量） & Hazeltine \& Meiss（回旋运动） \\
第 4 章 & 流体描述基础 & 第 3 章 & 第 4 章（Chapman--Enskog、Braginskii） & Bittencourt（输运推导） \\
第 5 章 & 磁流体力学 & 第 3 章、第 6 章 & 第 8 章（磁压、冻结、MHD 波） & Freidberg \textit{Ideal MHD} \\
第 6 章 & 等离子体中的波 & 第 4 章 & 第 5、6、7 章（冷、非均匀、温等离子体） & Stix \textit{Waves in Plasmas} \\
第 7 章 & 动理学理论 & 第 7 章 & 第 3 章、第 7 章（Landau 阻尼、速度空间不稳定性） & Krall \& Trivelpiece \\
第 8 章 & 碰撞与输运 & 第 5 章 & 第 3 章、第 4 章（库仑对数、Rosenbluth 势） & Helander \& Sigmar（新经典输运） \\
第 9 章 & 等离子体不稳定性 & 第 6 章、第 8 章 & 第 7 章、第 9 章（撕裂模、重联） & Freidberg；Wesson \textit{Tokamaks} \\
第 10 章 & 约束与受控核聚变 & 第 9、10 章 & 第 8 章（太阳风、激波） & Wesson \textit{Tokamaks}；Goldston \& Rutherford \\
第 11 章 & 低温工业等离子体 & 第 10 章 & 第 4 章末（Langmuir 鞘层与探针） & Lieberman \& Lichtenberg \\
第 12 章 & 等离子体基础诊断 & 第 10 章 & 第 4 章（探针理论） & Hutchinson \textit{Principles of Plasma Diagnostics} \\
\end{longtable}
\normalsize


\section{每章“外文教材独有视角”}

下面每条是\textbf{本书有意没展开、但在上述教材中讲得更透}的点，也是“读第二本书”时最该先看的地方。

\begin{enumerate}
\item \textbf{第 1 章}：Chen 把单位制与矢量恒等式集中放进附录 A，Fitzpatrick 则把量纲分析直接融入参数定义。两种做法都值得学：前者便于速查，后者让量纲成为物理直觉的一部分。
\item \textbf{第 2 章}：Fitzpatrick 第 1 章用“德拜屏蔽 $\to$ 等离子体参数 $\to$ 碰撞 $\to$ 磁化”的顺序定义等离子体，并配太阳风、电离层的实测数字。比纯理论定义更直观。
\item \textbf{第 3 章}：Fitzpatrick 第 2 章引入“平均化方法”与 Poincaré 不变量的讨论，解释了为什么绝热不变量在长时间尺度上守恒。本书只给结论。
\item \textbf{第 4 章}：Fitzpatrick 第 4 章系统给出 Chapman--Enskog 展开与 Braginskii 输运系数，是本书第 4、8 章之间缺失的那一环。
\item \textbf{第 5 章}：Fitzpatrick 第 8 章用太阳风与 Parker 螺旋、Cowling 反发电机定理把 MHD 放到天体物理语境中。动机比纯聚变视角更宽。
\item \textbf{第 6 章}：Fitzpatrick 把波分成三章（冷、非均匀、温等离子体）。其中非均匀等离子体的 WKB 处理与共振层、射线追踪，是本书完全没有覆盖的内容。
\item \textbf{第 7 章}：Fitzpatrick 第 7 章给出等离子体色散函数 $Z(\zeta)$ 的完整性质表与速度空间不稳定性的统一处理。本书只讨论了 Landau 阻尼的主线。
\item \textbf{第 8 章}：Fitzpatrick 第 3 章的 Landau 碰撞算符与 Rosenbluth 势比本书形式化得多。Chen 第 5 章则给出与实验对照的扩散、电阻率讨论。一个补数学，一个补实验。
\item \textbf{第 9 章}：Fitzpatrick 第 9 章用渐近匹配（asymptotic matching）求解撕裂模，并讨论 Sweet--Parker 与 plasmoid 不稳定性。这是理解现代重联研究的标准入口。
\item \textbf{第 10 章}：Goldston \& Rutherford 与 Wesson 提供真实装置尺度的数值（密度、温度、约束时间、阈值），可以把本书的劳森判据估算校准到工程现实。
\item \textbf{第 11 章}：Fitzpatrick 第 4 章末把鞘层与探针作为流体理论的直接应用；Lieberman \& Lichtenberg 则从刻蚀、沉积等工艺角度组织内容。这是“低温等离子体工程”的标准参考。
\item \textbf{第 12 章}：Hutchinson 按“测什么物理量”（密度、温度、磁场、电流分布）组织诊断。与本书按技术（探针、光谱、干涉、汤姆孙散射）组织的方式互为补充，适合做诊断选型的依据。
\end{enumerate}


\section{进阶专著速查}

{\small
\begin{longtable}{@{}p{4.4cm}p{5.2cm}p{4.9cm}@{}}
\caption{进入专门方向时的参考书}\label{tab:app-j-advanced}\\
\toprule
专著 & 出版信息 & 什么时候读它 \\
\midrule
\endfirsthead
\toprule
专著 & 出版信息 & 什么时候读它 \\
\midrule
\endhead
\midrule
\multicolumn{3}{r}{\small（续下页）}\\
\endfoot
\bottomrule
\endlastfoot
Wesson, \textit{Tokamaks} & Oxford University Press（初版 1987） & 处理托卡马克的具体不稳定性、约束与运行极限时 \\
Krall \& Trivelpiece, \textit{Principles of Plasma Physics} & McGraw-Hill, 1973 & 需要最完整的动理学与波的数学处理时 \\
Hazeltine \& Meiss, \textit{Plasma Confinement} & Addison-Wesley, 1992 & 研究约束、回旋动理学与输运理论时 \\
Freidberg, \textit{Ideal MHD} & Cambridge University Press, 2014 & 系统学习 MHD 平衡与稳定性、能量原理时 \\
Goedbloed \& Poedts, \textit{Principles of Magnetohydrodynamics} & Cambridge University Press, 2004 & 需要 MHD 谱理论（连续谱、阿尔芬本征模）时 \\
Bittencourt, \textit{Fundamentals of Plasma Physics} & Pergamon, 1986；Springer 第 3 版, 2004 & 需要比本书更详细的输运与波推导时 \\
Hutchinson, \textit{Principles of Plasma Diagnostics} & Cambridge University Press（初版 1987；第 2 版 2002） & 设计或解释诊断测量时 \\
Helander \& Sigmar, \textit{Collisional Transport in Magnetized Plasmas} & Cambridge University Press, 2002 & 研究新经典输运与杂质输运时 \\
Stix, \textit{Waves in Plasmas} & American Institute of Physics, 1992 & 需要波支分类与射频加热的系统处理时 \\
Bellan, \textit{Fundamentals of Plasma Physics} & Cambridge University Press, 2006 & 需要从第一性原理严格推导基础方程时 \\
\end{longtable}
}


\section{免费在线资源}

\begin{itemize}
\item \textbf{Fitzpatrick 教材书页}：\url{https://farside.ph.utexas.edu/Books/Plasma2/Plasma2.html}\\
作者提供的第 2 版在线版本（HTML），与纸质书内容对应，适合对照本附录中的章号。
\item \textbf{Fitzpatrick 等离子体物理讲义}\\
\url{https://farside.ph.utexas.edu/teaching/plasma/Plasmahtml/Plasmahtml.html}\\
单学期课程的讲义，篇幅更短，适合作为读本书时的并行材料。
\item \textbf{粒子模拟示例程序}：\url{https://farside.ph.utexas.edu/teaching/plasma/pic.tgz}\\
作者配套的 PIC 程序（含双流不稳定性算例），可用于验证本书第 9 章的线性增长率。
\end{itemize}

上述三个链接均为作者个人主页公开资源，核对日期见下节。


\section{元数据与核对说明}

本附录的章号与出版信息按以下方式核对（核对日期：2026-10-11）：

\begin{itemize}
\item Chen 第 3 版目录：经图书馆编目与出版方书目信息交叉核对，确认第 3 版\textbf{没有}单独的“碰撞效应”章。碰撞内容散见于第 5 章（扩散与电阻率）与第 7 章（动理学理论）。
\item Fitzpatrick 第 2 版目录：取自作者官方书页（本附录“免费在线资源”第一条），逐章逐节核对，共 9 章。
\item Goldston \& Rutherford：只核实到“六大部分”的主题划分，因此本附录对其采用\textbf{主题级}定位，不给章号。
\item 进阶专著的出版社、版次与年份：经开放图书馆（OpenLibrary）书目数据核对。
\item 本附录不引用任何未获授权的电子副本；Fitzpatrick 的在线版本为作者本人公开提供。
\end{itemize}

\begin{tip}{对照表用法}{tip-appj-howto}
\textbf{怎么看这张对照表：}不要按“章号对齐”去读书，而要按\textbf{问题}去读。当你卡在某个推导时，先在表中找到本书章节所在的行，再看哪一列的主题名与你的问题最接近，然后去那一章找对应的主题名。
\end{tip}
